\section{Stress, anxiety and the prefrontal executive system}\label{sec:stress}
\subsection{Mechanistic neuroscience}
The most developed account is due to Arnsten. Her review of the molecular signalling underlying stress effects describes the prefrontal cortex (PFC) as the most evolved brain region and also the one most sensitive to the detrimental effects of stress exposure \citep{arnsten2009stress}. The proposed mechanism is an inverted-U in catecholamine signalling: moderate norepinephrine and dopamine levels support prefrontal network firing and working memory, while the high levels of uncontrollable stress weaken them, and strengthen amygdala and striatal habit circuitry. A later review summarises the result as stress rapidly flipping the brain from reflective to reflexive control of behaviour \citep{arnsten2015stress}. Evolutionarily this is intelligible: a threat favours fast reflexes over deliberation.

\emph{Caveats.} Much of the mechanistic evidence comes from rodents and non-human primates. ``Offline'' is a popularisation of \emph{impaired}: the data support graded, dose-dependent impairment, not an on/off switch. Mild or controllable stress can leave executive function intact or improved. \citet{brewer2021unwinding} and \citet{brewer2016habit} state the strong version for a lay audience (anxiety shuts the PFC down and, with it, willpower and reasoning). We read these as clinician assertions that the primary literature supports only in graded form.
\citegap{human acute-stress meta-analyses of working memory, inhibition and shifting; reappraisal under stress}

\subsection{Cognitive psychology of anxiety}
Attentional Control Theory proposes that anxiety impairs the inhibition and shifting functions of executive control, so that anxious individuals are more captured by threat and worse at redirecting attention \citep{eysenck2007act}. An important nuance is that the theory predicts \emph{efficiency} (effort and time) to suffer more than \emph{effectiveness} (accuracy), because anxious people compensate with extra effort. That compensation is expensive for exactly the kind of person who has little executive reserve, which is the point of the next section.
